% Mathematical Interpretation of Vision-Based UAV Target Tracking
\documentclass[12pt,a4paper]{article}

\usepackage{amsmath,amssymb,bm}
\usepackage{geometry}
\usepackage{booktabs}
\usepackage{array}

\geometry{margin=1in}

\title{\textbf{Mathematical Interpretation of Vision-Based UAV Target Tracking}}
\author{}
\date{}

\begin{document}
\maketitle

\section{Overview}

This project implements a vision-based moving-target tracking system for a quadrotor UAV in a PyBullet simulation. A colored target moves through the environment and is observed using the UAV's camera. Computer vision detects the target in the image, and a feedback controller uses the detected image position and apparent target size to control the UAV.

The complete mathematical operation is

\begin{equation}
\boxed{
\text{Target Motion}
\rightarrow
\text{Camera Observation}
\rightarrow
\text{Target Detection}
\rightarrow
\text{Tracking Error}
\rightarrow
\text{PD Control}
\rightarrow
\text{UAV Motion}
\rightarrow
\text{New Observation}
}
\end{equation}

Thus, the project is a closed-loop visual feedback control system.

\section{1. Moving Target Model}

The target position in world coordinates is

\begin{equation}
\mathbf{p}_T(t)=
\begin{bmatrix}
x_T(t)\\y_T(t)\\z_T(t)
\end{bmatrix}.
\end{equation}

For the default circular trajectory used in the code,

\begin{equation}
x_T(t)=x_C+R\cos(\omega t),
\end{equation}

\begin{equation}
y_T(t)=y_C+R\sin(\omega t),
\end{equation}

\begin{equation}
z_T(t)=h.
\end{equation}

Therefore,

\begin{equation}
\boxed{
\mathbf{p}_T(t)=
\begin{bmatrix}
x_C+R\cos(\omega t)\\
y_C+R\sin(\omega t)\\
h
\end{bmatrix}
}
\end{equation}

where $R$ is the trajectory radius, $\omega$ is the speed parameter, $(x_C,y_C)$ is the trajectory centre, and $h$ is the target altitude.

For the current implementation,

\begin{equation}
R=1.5\,\mathrm{m},\qquad
\omega=0.4,\qquad
h=1.0\,\mathrm{m}.
\end{equation}

The target velocity is

\begin{equation}
\mathbf{v}_T(t)=
\frac{d\mathbf{p}_T(t)}{dt}
=
\begin{bmatrix}
-R\omega\sin(\omega t)\\
R\omega\cos(\omega t)\\
0
\end{bmatrix}.
\end{equation}

\section{2. UAV--Target Geometry}

Let the UAV position be

\begin{equation}
\mathbf{p}_D(t)=
\begin{bmatrix}
x_D(t)\\y_D(t)\\z_D(t)
\end{bmatrix}.
\end{equation}

The relative target position is

\begin{equation}
\boxed{
\mathbf{p}_{rel}(t)=\mathbf{p}_T(t)-\mathbf{p}_D(t)
}
\end{equation}

and the relative distance is

\begin{equation}
d(t)=\|\mathbf{p}_{rel}(t)\|.
\end{equation}

The controller does not directly use this 3-D distance. Instead, it estimates relative distance indirectly from the target's apparent image size.

\section{3. Camera Projection}

Let the target position in the camera coordinate system be

\begin{equation}
\mathbf{p}_{TC}=
\begin{bmatrix}
X_C\\Y_C\\Z_C
\end{bmatrix}.
\end{equation}

A perspective camera projects this 3-D point onto image coordinates:

\begin{equation}
\boxed{
u=f_x\frac{X_C}{Z_C}+c_x
}
\end{equation}

and

\begin{equation}
\boxed{
v=f_y\frac{Y_C}{Z_C}+c_y.
}
\end{equation}

Here $(u,v)$ are pixel coordinates, $f_x,f_y$ are focal lengths, and $(c_x,c_y)$ is the principal point.

PyBullet performs the actual rendering in the implementation, but these equations explain mathematically why changes in UAV position and orientation change the target's image position.

\section{4. HSV Color Segmentation}

The RGB camera image is converted to HSV:

\begin{equation}
I_{RGB}(u,v)\rightarrow I_{HSV}(u,v).
\end{equation}

The default green-target threshold is

\begin{equation}
35\leq H\leq85,
\end{equation}

\begin{equation}
80\leq S\leq255,
\end{equation}

\begin{equation}
80\leq V\leq255.
\end{equation}

The binary mask is therefore

\begin{equation}
\boxed{
M(u,v)=
\begin{cases}
1,&35\leq H\leq85,\;80\leq S\leq255,\;80\leq V\leq255,\\
0,&\text{otherwise}.
\end{cases}
}
\end{equation}

Morphological erosion and dilation are then applied to reduce small noisy regions.

\section{5. Target Detection}

Contours are extracted from the binary mask. The largest valid contour is selected:

\begin{equation}
\boxed{
C^*=\underset{C_i}{\operatorname{argmax}}\;A(C_i)
}
\end{equation}

subject to

\begin{equation}
A(C_i)\geq15.
\end{equation}

A minimum enclosing circle is fitted to the selected contour. The visual measurement returned by the tracker is

\begin{equation}
\boxed{
\mathbf{s}=
\begin{bmatrix}
c_x\\c_y\\r_{px}
\end{bmatrix}
}
\end{equation}

where $c_x,c_y$ are the target centre coordinates in pixels and $r_{px}$ is its apparent radius.

\section{6. Desired Visual State}

The camera resolution is

\begin{equation}
W=320,\qquad H=240.
\end{equation}

Hence the desired image centre is

\begin{equation}
c_x^*=\frac{W}{2}=160,
\end{equation}

\begin{equation}
c_y^*=\frac{H}{2}=120.
\end{equation}

The controller uses a desired target radius of

\begin{equation}
r_{px}^*=45.
\end{equation}

Therefore,

\begin{equation}
\boxed{
\mathbf{s}^*=
\begin{bmatrix}
160\\120\\45
\end{bmatrix}.
}
\end{equation}

The overall visual objective is

\begin{equation}
\boxed{
\mathbf{s}\rightarrow\mathbf{s}^*
}
\end{equation}

or

\begin{equation}
\boxed{
c_x\rightarrow160,\qquad
c_y\rightarrow120,\qquad
r_{px}\rightarrow45.
}
\end{equation}

\section{7. Visual Tracking Errors}

The horizontal error is

\begin{equation}
\boxed{
e_x=c_x-160.
}
\end{equation}

The vertical error is

\begin{equation}
\boxed{
e_z=120-c_y.
}
\end{equation}

The distance-related error is

\begin{equation}
\boxed{
e_d=45-r_{px}.
}
\end{equation}

Thus, the visual error vector can be written as

\begin{equation}
\boxed{
\mathbf{e}=
\begin{bmatrix}
e_x\\e_z\\e_d
\end{bmatrix}.
}
\end{equation}

Successful tracking corresponds to

\begin{equation}
\boxed{
\mathbf{e}\rightarrow\mathbf{0}.
}
\end{equation}

\section{8. Apparent Target Size and Distance}

The project uses the apparent radius of the target as an indirect distance measurement.

For a target of physical radius $R_T$, perspective projection gives approximately

\begin{equation}
\boxed{
r_{px}\approx\frac{fR_T}{Z}
}
\end{equation}

where $f$ is the camera focal length and $Z$ is target depth from the camera.

Therefore,

\begin{equation}
r_{px}\propto\frac{1}{Z}.
\end{equation}

Hence,

\begin{equation}
r_{px}\downarrow\Rightarrow Z\uparrow
\end{equation}

and

\begin{equation}
r_{px}\uparrow\Rightarrow Z\downarrow.
\end{equation}

This explains why the controller can use target size to determine whether the UAV should move forward or backward, even though the code does not explicitly calculate $Z$.

\section{9. PD Control Law}

The class in the code is named \texttt{PIDController}, but the integral gain is set to zero:

\begin{equation}
K_I=0.
\end{equation}

Therefore, the implemented controller is mathematically a PD controller.

The discrete-time control law is

\begin{equation}
\boxed{
u(k)=K_Pe(k)+
K_D\frac{e(k)-e(k-1)}{\Delta t}.
}
\end{equation}

The control frequency is

\begin{equation}
f_{ctrl}=48\,\mathrm{Hz},
\end{equation}

so

\begin{equation}
\boxed{
\Delta t=\frac{1}{48}\approx0.02083\,\mathrm{s}.
}
\end{equation}

The proportional term responds to the current error, while the derivative term responds to the rate of change of the error and helps reduce oscillation.

\section{10. Yaw Control}

The horizontal error controls the UAV yaw rate.

The implemented gains are

\begin{equation}
K_{P,yaw}=0.010,\qquad
K_{D,yaw}=0.002.
\end{equation}

Therefore,

\begin{equation}
\boxed{
\dot{\psi}
=
\operatorname{sat}_{[-2,2]}
\left[
0.010e_x+
0.002
\frac{e_x(k)-e_x(k-1)}
{\Delta t}
\right].
}
\end{equation}

The objective is

\begin{equation}
c_x\rightarrow160.
\end{equation}

\section{11. Altitude Control}

The vertical error controls UAV vertical velocity.

The gains are

\begin{equation}
K_{P,z}=0.008,\qquad
K_{D,z}=0.002.
\end{equation}

Thus,

\begin{equation}
\boxed{
v_z=
\operatorname{sat}_{[-0.8,0.8]}
\left[
0.008e_z+
0.002
\frac{e_z(k)-e_z(k-1)}
{\Delta t}
\right].
}
\end{equation}

The objective is

\begin{equation}
c_y\rightarrow120.
\end{equation}

\section{12. Forward/Backward Control}

The distance-related error controls forward and backward velocity.

The gains are

\begin{equation}
K_{P,d}=0.05,\qquad
K_{D,d}=0.01.
\end{equation}

The command is

\begin{equation}
\boxed{
v_x=
\operatorname{sat}_{[-1.5,1.5]}
\left[
0.05e_d+
0.01
\frac{e_d(k)-e_d(k-1)}
{\Delta t}
\right].
}
\end{equation}

If $r_{px}<45$, then $e_d>0$ and the controller commands forward motion.

If $r_{px}>45$, then $e_d<0$ and the controller commands backward motion.

Thus,

\begin{equation}
r_{px}\rightarrow45
\end{equation}

represents maintaining the desired apparent distance.

\section{13. Body-to-World Velocity Transformation}

The forward command $v_x$ is defined relative to the UAV's current heading.

Let the UAV rotation matrix be $\mathbf{R}_D$. Its forward direction in world coordinates is

\begin{equation}
\boxed{
\mathbf{f}_D=
\mathbf{R}_D
\begin{bmatrix}
1\\0\\0
\end{bmatrix}.
}
\end{equation}

The commanded world-frame velocity is

\begin{equation}
\boxed{
\mathbf{v}_{world}
=
v_x\mathbf{f}_D+
\begin{bmatrix}
0\\0\\v_z
\end{bmatrix}.
}
\end{equation}

This transformation is essential because the UAV's forward direction changes when its yaw changes.

\section{14. Desired Position and Yaw}

The tracking velocity is converted into a short-term desired position using discrete integration:

\begin{equation}
\boxed{
\mathbf{p}_{cmd}(k+1)
=
\mathbf{p}_D(k)+
\mathbf{v}_{world}(k)\Delta t.
}
\end{equation}

The commanded yaw is updated by integrating the yaw rate:

\begin{equation}
\boxed{
\psi_{cmd}(k+1)
=
\psi_{cmd}(k)+
\dot{\psi}(k)\Delta t.
}
\end{equation}

The desired attitude supplied to the low-level controller is

\begin{equation}
\boxed{
\boldsymbol{\phi}_{cmd}
=
\begin{bmatrix}
0\\0\\\psi_{cmd}
\end{bmatrix}.
}
\end{equation}

The low-level \texttt{DSLPIDControl} controller then converts the desired position and orientation into the drone's motor commands.

\section{15. Searching and Tracking}

The system has two states:

\begin{equation}
\boxed{
\text{SEARCHING}
\qquad\text{and}\qquad
\text{TRACKING}.
}
\end{equation}

When the target is detected,

\begin{equation}
\text{found}=1,
\end{equation}

the controller operates in TRACKING mode.

When the target is not detected,

\begin{equation}
\text{found}=0,
\end{equation}

the controller enters SEARCHING mode and commands

\begin{equation}
\boxed{
v_x=0,\qquad
v_y=0,\qquad
v_z=0,\qquad
\dot{\psi}=0.3.
}
\end{equation}

This produces a slow yaw rotation to reacquire the target.

\section{16. Closed-Loop System}

The complete system can be represented mathematically by

\begin{equation}
\mathbf{s}(k)
=
h\left(
\mathbf{x}_D(k),
\mathbf{p}_T(k)
\right),
\end{equation}

where $h(\cdot)$ represents the camera and target-detection process.

The visual error is

\begin{equation}
\mathbf{e}(k)
=
\mathbf{s}(k)-\mathbf{s}^*.
\end{equation}

The controller produces

\begin{equation}
\mathbf{u}(k)
=
\begin{bmatrix}
v_x(k)\\
v_z(k)\\
\dot{\psi}(k)
\end{bmatrix}.
\end{equation}

The UAV state evolves according to

\begin{equation}
\mathbf{x}_D(k+1)
=
f\left(
\mathbf{x}_D(k),
\mathbf{u}(k)
\right).
\end{equation}

The new UAV state produces a new camera observation:

\begin{equation}
\mathbf{s}(k+1)
=
h\left(
\mathbf{x}_D(k+1),
\mathbf{p}_T(k+1)
\right).
\end{equation}

Therefore, the closed-loop operation is

\begin{equation}
\boxed{
\mathbf{s}(k)
\rightarrow
\mathbf{e}(k)
\rightarrow
\mathbf{u}(k)
\rightarrow
\mathbf{x}_D(k+1)
\rightarrow
\mathbf{s}(k+1).
}
\end{equation}

This loop is repeated continuously throughout the simulation.

\section{17. Overall Tracking Objective}

A mathematical cost function that describes the tracking objective is

\begin{equation}
\boxed{
J=w_xe_x^2+w_ze_z^2+w_de_d^2
}
\end{equation}

where $w_x,w_z,w_d>0$ are weighting factors.

The ideal tracking condition is

\begin{equation}
e_x=0,\qquad
e_z=0,\qquad
e_d=0.
\end{equation}

Equivalently,

\begin{equation}
\boxed{
c_x=160,\qquad
c_y=120,\qquad
r_{px}=45.
}
\end{equation}

The implementation does not explicitly optimize $J$; this equation mathematically describes the desired behaviour of the controller.

\section{18. Overall Mathematical Interpretation}

The entire project can be summarized as follows:

\begin{equation}
\boxed{
\text{Moving Target}
\rightarrow
\text{3-D Camera Observation}
\rightarrow
\text{HSV Segmentation}
\rightarrow
(c_x,c_y,r_{px})
\rightarrow
(e_x,e_z,e_d)
\rightarrow
\text{PD Control}
\rightarrow
(v_x,v_z,\dot{\psi})
\rightarrow
\text{Coordinate Transformation}
\rightarrow
\text{Desired UAV State}
\rightarrow
\text{UAV Motion}
\rightarrow
\text{New Camera Observation}.
}
\end{equation}

The fundamental tracking condition is

\begin{equation}
\boxed{
\lim_{k\rightarrow\infty}
\begin{bmatrix}
c_x(k)-160\\
c_y(k)-120\\
r_{px}(k)-45
\end{bmatrix}
=
\mathbf{0}.
}
\end{equation}

In practical terms, the UAV continuously adjusts its yaw, altitude, and forward/backward movement so that the moving target remains centred in the camera image and maintains the desired apparent size.

\section{19. Key Mathematical Relationships}

The most important relationships in the project are:

\begin{equation}
\boxed{
c_x\rightarrow160
\quad\Rightarrow\quad
\text{horizontal alignment}
}
\end{equation}

\begin{equation}
\boxed{
c_y\rightarrow120
\quad\Rightarrow\quad
\text{vertical alignment}
}
\end{equation}

\begin{equation}
\boxed{
r_{px}\rightarrow45
\quad\Rightarrow\quad
\text{desired relative distance}
}
\end{equation}

and

\begin{equation}
\boxed{
u=K_Pe+K_D\frac{\Delta e}{\Delta t}
\quad\Rightarrow\quad
\text{corrective UAV motion}.
}
\end{equation}

Together, these equations explain the mathematical working of the vision-based UAV tracking system.

\end{document}
